\documentclass[11pt]{article}

\usepackage[utf8]{inputenc}
\usepackage[T1]{fontenc}
\usepackage[spanish,es-noshorthands]{babel}
\usepackage[margin=2.5cm]{geometry}
\usepackage{booktabs}
\usepackage{graphicx}
\usepackage{amsmath}
\usepackage{caption}
\usepackage{float}
\usepackage{xcolor}
\usepackage{url}
\usepackage{hyperref}
\hypersetup{colorlinks=true, linkcolor=blue!50!black, citecolor=black, urlcolor=blue!50!black}

\title{\textbf{Golden Hour}\\[4pt]
  \large Accesibilidad a Salud Resolutiva en el Perú:\\
  Un Análisis Geoespacial de Ica, Cusco y Loreto}
\author{Proyecto Integrador Golden Hour}
\date{\today}

\begin{document}

\maketitle

\renewcommand{\abstractname}{Resumen Ejecutivo}
\begin{abstract}
\noindent
El acceso oportuno a un establecimiento de salud con capacidad resolutiva
(niveles II y III, MINSA) es determinante para la supervivencia en
emergencias obstétricas, traumatológicas y quirúrgicas. Este estudio
construye y evalúa un indicador de accesibilidad geoespacial --el
tiempo de viaje hasta el establecimiento resolutivo más cercano-- para
tres departamentos que representan las tres macrorregiones naturales del
Perú: Ica (Costa), Cusco (Sierra) y Loreto (Selva). A partir de 8{,}632
centros poblados (fuente INEI, censo 1999/2002) y 62 establecimientos
resolutivos (RENIPRESS/SUSALUD), se calculan matrices origen-destino por
carretera, caminata y bicicleta, y se agregan en indicadores ponderados
por población a nivel distrital, provincial y departamental.

Los hallazgos son consistentes con una brecha de accesibilidad severa y
desigualmente distribuida. El coeficiente de Gini del tiempo de acceso
ponderado por población es de \textbf{0.729} a nivel nacional (de los 3
departamentos), con una dispersión notable entre ámbitos (Ica 0.768,
Loreto 0.655, Cusco 0.510). Bajo el umbral de 60 minutos configurado como
referencia de ``hora dorada'', aproximadamente \textbf{29.1\,\%} de la
población analizada queda fuera de ese margen. La diferencia entre
entornos urbano y rural es estadísticamente significativa (Mann-Whitney
$U$, $p = 0.0002$), con un tiempo de acceso ponderado casi 4 veces mayor
en zonas rurales (120.6 min) que en zonas urbanas (30.2 min). Un cruce
analítico exploratorio muestra una asociación positiva, moderada y
estadísticamente significativa entre el índice de ruralidad distrital y
el tiempo de acceso (Spearman $\rho = 0.266$, $p = 0.0001$, $n = 200$
distritos) -- una asociación espacial/correlacional, no causal. Se
discuten las limitaciones metodológicas del estudio (vintage censal,
disponibilidad de motores de ruteo, redes fluviales no mapeadas) y se
proponen recomendaciones de política focalizadas en los distritos
identificados como más críticos.
\end{abstract}

\tableofcontents

% =============================================================================
\section{Introducción y Contexto Geográfico}
% =============================================================================

La accesibilidad geográfica a servicios de salud resolutivos es una de las
dimensiones más citadas -- y más difíciles de medir con datos abiertos -- de
la equidad sanitaria. El concepto de \textit{golden hour} (``hora dorada''),
originado en la medicina de trauma, sostiene que la probabilidad de
supervivencia ante una emergencia grave cae abruptamente cuando el tiempo
transcurrido hasta recibir atención definitiva supera aproximadamente 60
minutos. En un país con la geografía tan heterogénea como el Perú --costa
desértica, cordillera andina y llanura amazónica-- ese umbral se cumple de
maneras radicalmente distintas según el territorio.

Este proyecto integrador selecciona un departamento por cada una de las
tres macrorregiones naturales del país para capturar esa heterogeneidad:

\begin{itemize}
  \item \textbf{Ica (Costa)}: territorio desértico costero, con una red vial
    relativamente densa y una población concentrada en pocos valles.
  \item \textbf{Cusco (Sierra)}: territorio andino de orografía compleja,
    con provincias altamente pobladas cerca de la capital regional y
    provincias remotas de selva alta (p.\,ej. La Convención) con muy baja
    densidad vial.
  \item \textbf{Loreto (Selva)}: el departamento más extenso del Perú,
    donde el transporte fluvial sustituye en gran medida a la red vial
    terrestre -- una limitación central que se discute en la
    Sección~\ref{sec:limitaciones}.
\end{itemize}

El objetivo del proyecto es construir, de extremo a extremo, un pipeline
reproducible --ingesta de datos abiertos, validación, cálculo de tiempos de
viaje, construcción de indicadores y un tablero interactivo-- que permita
cuantificar la brecha de acceso a salud resolutiva en estos tres territorios
y identificar, con evidencia, dónde una intervención de infraestructura vial
o de red de referencia tendría el mayor impacto marginal.

% =============================================================================
\section{Metodología y Decisiones Técnicas}
% =============================================================================

El proyecto se organizó en cuatro fases de código (\texttt{acquisition.py},
\texttt{validation.py}, \texttt{routing.py}, \texttt{metrics.py} /
\texttt{export.py}) más un tablero interactivo (\texttt{dashboard.py}), todas
parametrizadas desde un único archivo de configuración
(\texttt{config.md}). Esta sección documenta las decisiones técnicas no
triviales tomadas en cada etapa.

\subsection{Ingesta y validación de datos}

\subsubsection{Fuentes y sustituciones de fuente}

Se utilizaron tres fuentes de datos abiertos del Estado peruano:

\begin{itemize}
  \item \textbf{RENIPRESS} (SUSALUD/MINSA): registro de establecimientos de
    salud, resuelto dinámicamente vía la API CKAN del Portal Nacional de
    Datos Abiertos para obtener siempre el archivo mensual más reciente.
  \item \textbf{Límites administrativos} (IGN): departamentos, provincias y
    distritos, en formato \textit{shapefile}.
  \item \textbf{Centros poblados}: usados como puntos de demanda.
\end{itemize}

La fuente de centros poblados requirió dos sustituciones sucesivas,
documentadas explícitamente en \texttt{config.md} como decisiones de
diseño:

\begin{enumerate}
  \item \textbf{MINEDU/SIGMED $\rightarrow$ IGN.} El visor SIGMED
    (\url{sigmed.minedu.gob.pe}) no expone una URL de descarga estática:
    el botón de descarga dispara un \textit{postback} ASP.NET parametrizado
    por selección interactiva de departamento/provincia/distrito, sin
    endpoint de archivo verificable programáticamente. Se sustituyó por el
    dataset equivalente ``Centros Poblados'' del IGN, publicado en el mismo
    portal de datos abiertos que RENIPRESS, con URL estable.
  \item \textbf{IGN $\rightarrow$ INEI (censo 1999/2002, vía INGEMMET/GEOCATMIN).}
    La validación de Fase 1 detectó que el 100\,\% de los registros de
    Loreto en el dataset del IGN tenían la columna de latitud duplicada con
    el valor de longitud (\texttt{Y = X}) -- un defecto sistemático de la
    fuente, no un error de intercambio de columnas corregible
    automáticamente-- dejando ese departamento sin puntos de demanda
    válidos. Adicionalmente, el dataset del IGN no incluye ningún campo de
    población, lo cual es incompatible con el requisito central de Fase 3
    de ponderar los indicadores por habitantes. Se sustituyó, para los 3
    departamentos de forma consistente, por la capa ``Centros Poblados -
    INEI'' publicada como servicio ArcGIS REST por INGEMMET/GEOCATMIN, que
    aporta población real (\texttt{TOT\_POB99}) y clasificación
    urbano/rural oficial (\texttt{CLASIF02}).
\end{enumerate}

\subsubsection{Obstáculos técnicos de adquisición}

La adquisición programática de datos abiertos del Estado peruano presentó
dos obstáculos no documentados en las especificaciones públicas de los
portales:

\begin{itemize}
  \item \textbf{WAF con bloqueo por \textit{User-Agent} (HTTP 418).} El
    portal \url{datosabiertos.gob.pe} está protegido por un
    \textit{Web Application Firewall} que devuelve el código HTTP 418
    (``I'm a teapot'') ante el \textit{User-Agent} por defecto de la
    librería \texttt{requests} de Python (\texttt{python-requests/x.y}),
    tanto para la API como para la descarga de archivos. Se resolvió
    enviando encabezados HTTP consistentes con un navegador real.
  \item \textbf{Forma de respuesta no estándar de la API CKAN.} La API CKAN
    de este portal envuelve el paquete solicitado (\texttt{package\_show})
    en una lista de un elemento en el campo \texttt{result}, en vez de
    devolver el objeto directamente como es usual en CKAN. El código de
    resolución dinámica de RENIPRESS soporta ambas formas defensivamente.
  \item \textbf{Certificado TLS incompleto en INGEMMET/GEOCATMIN.} El
    servidor \url{geocatmin.ingemmet.gob.pe} no envía la cadena de
    certificados intermedia completa, lo que provoca un error de
    verificación TLS en cualquier cliente HTTPS estándar (no solo en este
    proyecto). Verificado manualmente que el contenido servido es
    correcto, se omite la verificación TLS únicamente para este host y
    únicamente para la lectura de datos públicos de solo lectura.
  \item \textbf{Límite de 1{,}000 registros sin paginación.} El servicio
    ArcGIS REST de centros poblados no soporta paginación por
    \textit{offset} y limita cada consulta a 1{,}000 registros. Se paginó
    manualmente consultando cada provincia por separado; la provincia de
    La Convención (Cusco) aun así excedió el límite y sus registros
    quedaron truncados a 1{,}000 (ver Sección~\ref{sec:limitaciones}).
\end{itemize}

\subsubsection{Reglas de validación}

Sobre ambas fuentes de puntos (establecimientos y centros poblados) se
aplicó una batería de reglas de calidad, documentada en su totalidad en
\texttt{logs/data\_quality\_report.csv}: coordenadas nulas, en cero o
faltantes (descartadas); coordenadas fuera del \textit{bounding box}
continental del Perú; latitud/longitud intercambiadas, con corrección
automática y tasa de recuperación explícita; puntos fuera de su polígono
distrital declarado (validación espacial vía \textit{spatial join},
conservados con bandera para revisión); códigos de establecimiento
duplicados; problemas de codificación UTF-8 vs. Latin-1; y normalización de
categorías resolutivas (II-1, II-2, II-E, III-1, III-2, III-E) frente a no
resolutivas. Tras la validación quedaron 2{,}312 establecimientos de salud
(62 resolutivos: 26 en Cusco, 21 en Ica, 15 en Loreto) y 8{,}632 centros
poblados de demanda (5{,}170 Cusco, 2{,}271 Loreto, 1{,}191 Ica), sobre 212
distritos y 24 provincias.

\subsection{Motor de ruteo, \textit{snapping} y \textit{fallback} Haversine}

Dado que el universo de demanda (8{,}632 puntos) supera el techo
configurado de 5{,}000 puntos, se aplicó un muestreo estratificado por
distrito, ponderado por población real (censo INEI), reduciendo la log de
la población para evitar que un puñado de asentamientos muy grandes
dominara por completo el muestreo ponderado sin reemplazo. El resultado fue
una muestra de 5{,}000 puntos (57.9\,\% del universo).

El motor de ruteo sigue una arquitectura de tolerancia a fallos en dos
niveles:

\begin{enumerate}
  \item \textbf{OSRM (Docker local) como motor primario}, únicamente para
    el perfil \textit{car}: el contenedor documentado en
    \texttt{config.md} expone un único perfil (\textit{driving}), por lo
    que no puede rutear caminata ni bicicleta de forma correcta.
  \item \textbf{OSMnx/NetworkX como \textit{fallback} automático}, para
    \textit{car} cuando OSRM no está disponible, y siempre para
    \textit{foot}/\textit{bike}. Cada punto se asocia (\textit{snapping})
    al nodo vial más cercano del grafo descargado de OpenStreetMap, y se
    calcula la ruta más rápida mediante el algoritmo de Dijkstra de una
    sola fuente por origen.
\end{enumerate}

En esta ejecución, OSRM no estuvo disponible (no hay contenedor Docker
activo en el entorno de cómputo), por lo que el cálculo recayó en el
\textit{fallback}. Se detectó empíricamente que, una vez que el muestreo
estratificado dispersa los puntos de demanda por múltiples distritos, el
\textit{bounding box} combinado por departamento excede ampliamente el área
recomendada para una única consulta a la API pública de Overpass (hasta 475
veces el área máxima recomendada para Loreto), generando cientos de
sub-consultas que agotan el tiempo de espera. Se implementó un umbral de
área práctica (5\,grados\textsuperscript{2}) por encima del cual el sistema
omite el intento de descarga del grafo y recurre directamente al
\textit{fallback} de distancia Haversine:
\[
  d_{\text{hav}}(i,j) = 2 R \arcsin\!\left(\sqrt{\sin^2\!\tfrac{\Delta\phi}{2} + \cos\phi_i \cos\phi_j \sin^2\!\tfrac{\Delta\lambda}{2}}\right),
  \qquad R = 6{,}371.0088\text{ km},
\]
multiplicada por un \textbf{factor de desvío empírico de 1.4} para
aproximar la distancia real de red a partir de la distancia en línea recta,
y convertida a tiempo con una velocidad asumida por modo (40, 5 y 15
km/h para \textit{car}, \textit{foot} y \textit{bike}, respectivamente).
Cada punto queda marcado con \texttt{is\_unroutable = True} cuando se usa
esta estimación. En la ejecución reportada, el 100\,\% de los pares
origen-destino recayeron en este \textit{fallback} para los 3
departamentos -- una limitación central discutida en la
Sección~\ref{sec:limitaciones}, y precisamente el escenario para el que
OSRM (que opera sobre extractos OSM locales sin las restricciones de la
API pública de Overpass) está documentado como motor primario.

\subsection{Formulación del coeficiente de Gini espacial y la curva de Lorenz}
\label{sec:gini-formula}

El coeficiente de Gini se calculó sobre el tiempo de acceso en automóvil
$t_i$ de cada punto de demanda $i$, ponderado por su población $w_i$. Sea
$(t_{(1)}, \dots, t_{(n)})$ el tiempo de acceso ordenado de forma
\emph{ascendente}, con $w_{(1)}, \dots, w_{(n)}$ las poblaciones asociadas.
Se definen las fracciones acumuladas de población y de ``tiempo total''
hasta el punto $k$:
\[
  p_k = \frac{\sum_{i=1}^{k} w_{(i)}}{\sum_{i=1}^{n} w_{(i)}},
  \qquad
  L_k = \frac{\sum_{i=1}^{k} w_{(i)}\, t_{(i)}}{\sum_{i=1}^{n} w_{(i)}\, t_{(i)}},
\]
con $p_0 = L_0 = 0$. La curva de Lorenz es la función lineal a trozos que
interpola los puntos $(p_k, L_k)$, y el coeficiente de Gini se obtiene como
$1$ menos el doble del área bajo esa curva (regla del trapecio):
\[
  G = 1 - 2\sum_{k=1}^{n} (p_k - p_{k-1}) \cdot \frac{L_k + L_{k-1}}{2}.
\]

\textbf{Justificación metodológica.} El Gini y la curva de Lorenz se usan
tradicionalmente para medir la desigualdad en la distribución de un
\emph{bien} (p.\,ej. el ingreso) entre personas: un Gini alto indica que
pocas personas concentran una porción desproporcionada del bien. Aquí se
aplican, de forma deliberada y no convencional, a un \emph{costo} (minutos
de viaje): un Gini alto no significa que pocos concentran un beneficio,
sino que la \emph{carga} del tiempo de traslado está muy desigualmente
repartida -- una fracción de la población soporta tiempos de acceso muchas
veces mayores que el resto. Esta adaptación es consistente con la
literatura de equidad en transporte y accesibilidad, que reporta
habitualmente un ``Gini de tiempos de viaje'' con esta misma lectura
invertida.

\subsection{Construcción de indicadores}

Todas las agregaciones jerárquicas (distrital, provincial, departamental)
usan como estadístico principal la \textbf{media ponderada por
población}, evitando explícitamente que un caserío de unas pocas decenas
de habitantes pese lo mismo que una capital de distrito de varios miles.
La media simple y la mediana se reportan solo como contraste diagnóstico.
El contraste urbano/rural usa la prueba de Mann-Whitney $U$ (no
paramétrica) en lugar de una prueba $t$, dado que los tiempos de viaje
exhiben colas largas y distribuciones marcadamente asimétricas que violan
el supuesto de normalidad de esta última.

\subsubsection{Cruce analítico: tiempo de acceso e índice de ruralidad}

Como segundo indicador regional para el cruce analítico exigido por la
Fase 3, se construyó un \textbf{índice de ruralidad distrital}: el
porcentaje de la población de cada distrito clasificada como ``RURAL''
según la variable oficial del INEI (\texttt{CLASIF02}), derivado
directamente de los datos ya validados, sin necesidad de una fuente
externa adicional. Este índice se correlacionó (Spearman, apropiado para
relaciones monótonas no necesariamente lineales) con el tiempo de acceso
ponderado a nivel distrital.

\textbf{Naturaleza de la relación (documentado también en el código y en
los logs de ejecución de \texttt{metrics.py}):} la correlación resultante
es una asociación \textbf{correlacional/espacial} entre unidades
geográficas agregadas (una correlación ecológica), \emph{no una relación
causal}. Un distrito más rural puede tener peor acceso por múltiples
factores confundidos entre sí --menor densidad vial, mayor distancia a la
capital provincial, orografía, dispersión poblacional-- que este cruce
bivariado no aísla ni controla.

% =============================================================================
\section{Resultados y Análisis Comparativo}
% =============================================================================

\subsection{Cobertura y tiempos de acceso}

La Tabla~\ref{tab:coverage-bands} muestra la distribución de la población
analizada por banda de tiempo de acceso en automóvil. Aunque solo el
19.8\,\% de los \emph{puntos} de demanda están a menos de 30 minutos de un
establecimiento resolutivo, esos puntos concentran el 57.6\,\% de la
\emph{población} -- evidencia directa de por qué la ponderación por
habitantes es indispensable: los asentamientos grandes tienden a ubicarse
cerca de la red de servicios, y un promedio simple por punto subestimaría
sistemáticamente el acceso real de la mayoría de las personas.

\begin{table}[htbp]
\centering
\caption{Cobertura por banda de tiempo de acceso a salud resolutiva}
\label{tab:coverage-bands}
\resizebox{\ifdim\width>\textwidth\textwidth\else\width\fi}{!}{%
\begin{tabular}{lrrrr}
\toprule
Banda de tiempo & Núm. centros poblados & Población & Pct. centros poblados & Pct. población \\
\midrule
<30 min & 988 & 1,023,536 & 19.76 & 57.62 \\
30-60 min & 1,119 & 235,719 & 22.38 & 13.27 \\
60-120 min & 1,625 & 277,158 & 32.50 & 15.60 \\
>120 min (fuera de Golden Hour) & 1,268 & 239,798 & 25.36 & 13.50 \\
\bottomrule
\end{tabular}%
}
\end{table}


Bajo el umbral de referencia de 60 minutos (\texttt{config.md}), el
29.1\,\% de la población analizada queda \textbf{fuera} de la ``hora
dorada'' vehicular (suma de las bandas 60--120 y $>$120 min).

La Tabla~\ref{tab:access-department} resume el tiempo de acceso ponderado
por departamento. La brecha es sustancial: Loreto (135.5 min) es más de
\textbf{10 veces} más lento que Ica (12.5 min), y más de el triple de
Cusco (46.4 min). Nótese también la magnitud de la diferencia entre la
media simple y la ponderada en Cusco (73.5 vs.\ 46.4 min) y en Loreto
(242.5 vs.\ 135.5 min): en ambos casos la población está sistemáticamente
mejor ubicada que el centro poblado promedio, lo que confirma que reportar
solo la media simple habría exagerado la severidad de la brecha para la
mayoría de los habitantes (sin restarle gravedad a la de los que sí están
mal servidos, como se ve en el análisis de desigualdad).

\begin{table}[htbp]
\centering
\caption{Tiempo de acceso ponderado por población, nivel departamental}
\label{tab:access-department}
\resizebox{\ifdim\width>\textwidth\textwidth\else\width\fi}{!}{%
\begin{tabular}{lrrrrr}
\toprule
Departamento & Núm. centros poblados & Población & Tiempo medio simple (min) & Mediana (min) & Tiempo ponderado (min) \\
\midrule
CUSCO & 2,992 & 734,566 & 73.48 & 65.39 & 46.38 \\
ICA & 692 & 531,187 & 39.66 & 29.29 & 12.51 \\
LORETO & 1,316 & 510,458 & 242.53 & 207.32 & 135.51 \\
\bottomrule
\end{tabular}%
}
\end{table}


\subsection{Brechas críticas por distrito}

La Tabla~\ref{tab:critical-districts} identifica los 10 distritos con peor
accesibilidad ponderada dentro de cada departamento. El caso más extremo es
\textbf{Yavarí} (provincia Mariscal Ramón Castilla, Loreto), con un tiempo
de acceso ponderado de 656.0 minutos (casi 11 horas) para una población de
2{,}298 habitantes -- consistente con la ausencia de red vial en esa zona
de frontera amazónica y la dependencia del transporte fluvial. En el otro
extremo, el peor distrito de Ica (Marcona, 134.8 min) tiene un tiempo de
acceso comparable al \emph{mejor} de los 10 distritos más críticos de
Loreto, ilustrando la magnitud de la brecha interregional.

\begin{table}[htbp]
\centering
\caption{Top 10 distritos con peor accesibilidad ponderada, por departamento}
\label{tab:critical-districts}
\resizebox{\ifdim\width>\textwidth\textwidth\else\width\fi}{!}{%
\begin{tabular}{lllrrrrrr}
\toprule
Departamento & Provincia & Distrito & Núm. centros poblados & Población & Tiempo medio simple (min) & Mediana (min) & Tiempo ponderado (min) & Ranking \\
\midrule
CUSCO & QUISPICANCHI & CAMANTI & 35 & 1,721 & 270.79 & 277.00 & 271.54 & 1 \\
CUSCO & QUISPICANCHI & MARCAPATA & 56 & 3,259 & 164.58 & 163.62 & 165.14 & 2 \\
CUSCO & PAUCARTAMBO & KOSÐIPATA & 32 & 2,622 & 157.80 & 160.20 & 163.61 & 3 \\
CUSCO & ESPINAR & CONDOROMA & 49 & 798 & 140.34 & 141.67 & 139.16 & 4 \\
CUSCO & ANTA & MOLLEPATA & 44 & 900 & 131.20 & 130.33 & 131.97 & 5 \\
CUSCO & QUISPICANCHI & OCONGATE & 22 & 6,645 & 123.32 & 122.12 & 123.83 & 6 \\
CUSCO & CALCA & YANATILE & 60 & 7,292 & 129.77 & 129.23 & 120.71 & 7 \\
CUSCO & QUISPICANCHI & CCARHUAYO & 21 & 2,327 & 120.20 & 122.36 & 114.17 & 8 \\
CUSCO & LA CONVENCION & VILCABAMBA & 120 & 7,021 & 100.28 & 86.17 & 109.10 & 9 \\
CUSCO & LA CONVENCION & ECHARATE & 143 & 27,412 & 104.80 & 89.89 & 108.50 & 10 \\
ICA & NAZCA & MARCONA & 5 & 17,299 & 114.60 & 120.61 & 134.81 & 1 \\
ICA & CHINCHA & SAN PEDRO DE HUACARPANA & 16 & 1,158 & 116.20 & 112.78 & 119.47 & 2 \\
ICA & PISCO & HUANCANO & 23 & 1,081 & 117.17 & 113.07 & 116.30 & 3 \\
ICA & PALPA & TIBILLO & 14 & 199 & 95.41 & 96.77 & 95.78 & 4 \\
ICA & CHINCHA & CHAVIN & 28 & 734 & 86.38 & 90.67 & 90.94 & 5 \\
ICA & CHINCHA & SAN JUAN DE YANAC & 19 & 648 & 75.95 & 79.01 & 82.61 & 6 \\
ICA & ICA & YAUCA DEL ROSARIO & 17 & 1,951 & 77.35 & 80.18 & 73.14 & 7 \\
ICA & ICA & OCUCAJE & 10 & 2,382 & 82.89 & 68.17 & 68.38 & 8 \\
ICA & PISCO & HUMAY & 21 & 4,597 & 60.67 & 59.44 & 54.61 & 9 \\
ICA & NAZCA & CHANGUILLO & 15 & 2,742 & 48.29 & 36.42 & 41.79 & 10 \\
LORETO & MARISCAL RAMON CASTILLA & YAVARI & 23 & 2,298 & 643.62 & 658.06 & 655.96 & 1 \\
LORETO & MARISCAL RAMON CASTILLA & RAMON CASTILLA & 43 & 10,946 & 635.93 & 628.65 & 651.30 & 2 \\
LORETO & ALTO AMAZONAS & MORONA & 28 & 5,126 & 539.64 & 528.39 & 509.26 & 3 \\
LORETO & MAYNAS & PUTUMAYO & 52 & 3,219 & 413.05 & 411.94 & 464.80 & 4 \\
LORETO & MARISCAL RAMON CASTILLA & SAN PABLO & 24 & 10,154 & 467.90 & 459.12 & 450.63 & 5 \\
LORETO & LORETO & TROMPETEROS & 18 & 3,564 & 454.59 & 483.22 & 432.43 & 6 \\
LORETO & ALTO AMAZONAS & MANSERICHE & 22 & 3,572 & 428.17 & 419.85 & 430.87 & 7 \\
LORETO & ALTO AMAZONAS & PASTAZA & 62 & 8,604 & 479.32 & 499.96 & 428.31 & 8 \\
LORETO & REQUENA & CAPELO & 6 & 3,376 & 412.86 & 407.74 & 418.97 & 9 \\
LORETO & REQUENA & TAPICHE & 6 & 1,550 & 405.74 & 398.95 & 417.55 & 10 \\
\bottomrule
\end{tabular}%
}
\end{table}


\subsection{Desigualdad en el acceso: Gini y curva de Lorenz}

La Figura~\ref{fig:lorenz} muestra la curva de Lorenz del tiempo de acceso
ponderado por población, junto con la línea de igualdad perfecta. El
coeficiente de Gini resultante (Tabla~\ref{tab:gini}) es de \textbf{0.729}
a nivel de los 3 departamentos combinados.

\begin{table}[htbp]
\centering
\caption{Coeficiente de Gini del tiempo de acceso, nacional y por departamento}
\label{tab:gini}
\resizebox{\ifdim\width>\textwidth\textwidth\else\width\fi}{!}{%
\begin{tabular}{lr}
\toprule
Ámbito & Gini \\
\midrule
NACIONAL (3 deptos) & 0.729 \\
CUSCO & 0.510 \\
ICA & 0.768 \\
LORETO & 0.655 \\
\bottomrule
\end{tabular}%
}
\end{table}


\begin{figure}[H]
  \centering
  \includegraphics[width=0.75\textwidth]{figures/lorenz_curve.pdf}
  \caption{Curva de Lorenz del tiempo de acceso a salud resolutiva, ponderada
    por población. El área entre la curva y la línea de igualdad perfecta
    es proporcional al coeficiente de Gini (Sección~\ref{sec:gini-formula}).}
  \label{fig:lorenz}
\end{figure}

Un hallazgo que solo emerge al desagregar por departamento (y no estaría
visible en el Gini nacional agregado) es que \textbf{Ica, pese a tener el
menor tiempo promedio, tiene la mayor desigualdad} (Gini = 0.768): la
mayoría de su población está muy bien servida, pero una minoría enfrenta
tiempos varias veces mayores, generando una distribución muy cóncava.
Cusco, en cambio, tiene la menor desigualdad (Gini = 0.510): su
accesibilidad es uniformemente intermedia, sin la misma polarización.
Loreto combina un tiempo promedio alto \emph{y} una desigualdad
considerable (Gini = 0.655) -- el peor perfil combinado de los tres.

\subsection{Distribución por departamento y contraste urbano/rural}

La Figura~\ref{fig:dist-dept} compara la distribución completa del tiempo
de acceso en automóvil entre los tres departamentos, con líneas de
referencia en 30, 60 y 120 minutos.

\begin{figure}[H]
  \centering
  \includegraphics[width=0.85\textwidth]{figures/access_time_distribution_by_department.pdf}
  \caption{Distribución del tiempo de acceso en automóvil por departamento
    (diagrama de caja, sin valores atípicos extremos para preservar la
    legibilidad de la escala).}
  \label{fig:dist-dept}
\end{figure}

El contraste urbano/rural (Tabla~\ref{tab:urban-rural}) es marcado: el
tiempo de acceso ponderado rural (120.6 min) es \textbf{casi 4 veces} el
urbano (30.2 min). La prueba de Mann-Whitney $U$ (Tabla~\ref{tab:urban-rural-test})
confirma que esta diferencia es estadísticamente significativa
($U = 438{,}052$, $p = 0.0002$), lo que descarta que la brecha observada se
deba al azar muestral.

\begin{table}[htbp]
\centering
\caption{Tiempo de acceso: contraste urbano vs. rural}
\label{tab:urban-rural}
\resizebox{\ifdim\width>\textwidth\textwidth\else\width\fi}{!}{%
\begin{tabular}{lrrrrr}
\toprule
Grupo & Núm. centros poblados & Población & Tiempo ponderado (min) & Tiempo medio simple (min) & Mediana (min) \\
\midrule
Urbano & 215 & 1,153,695 & 30.18 & 100.57 & 58.75 \\
Rural & 4,785 & 622,516 & 120.59 & 113.87 & 72.55 \\
\bottomrule
\end{tabular}%
}
\end{table}

\begin{table}[htbp]
\centering
\caption{Prueba de Mann-Whitney U: tiempo de acceso urbano vs. rural}
\label{tab:urban-rural-test}
\resizebox{\ifdim\width>\textwidth\textwidth\else\width\fi}{!}{%
\begin{tabular}{lrr}
\toprule
Método & Estadístico & p-valor \\
\midrule
Mann-Whitney U & 438,052 & 0.0002 \\
\bottomrule
\end{tabular}%
}
\end{table}


\subsection{Cruce con el índice de ruralidad distrital}

La correlación de Spearman entre el índice de ruralidad distrital y el
tiempo de acceso ponderado (Tabla~\ref{tab:rurality-correlation}) es
positiva, moderada y estadísticamente significativa
($\rho = 0.266$, $p = 0.0001$, $n = 200$ distritos): a mayor proporción de
población rural en un distrito, mayor tiende a ser su tiempo de acceso.
Como se documentó en la Sección~\ref{sec:gini-formula}, esta es una
asociación \textbf{correlacional/espacial}, no causal: el coeficiente es
moderado (no fuerte) precisamente porque factores no controlados en este
cruce bivariado --orografía, densidad vial, distancia a la capital
provincial-- explican una parte sustancial de la varianza restante.

\begin{table}[htbp]
\centering
\caption{Correlación (Spearman) entre índice de ruralidad y tiempo de acceso distrital}
\label{tab:rurality-correlation}
\resizebox{\ifdim\width>\textwidth\textwidth\else\width\fi}{!}{%
\begin{tabular}{llrrr}
\toprule
Método & Naturaleza & Coeficiente & p-valor & Núm. distritos \\
\midrule
Spearman & correlacional/espacial (no causal) & 0.266 & 0.0001 & 200 \\
\bottomrule
\end{tabular}%
}
\end{table}


\subsection{Mapa de brechas distritales}

La Figura~\ref{fig:gap-map} presenta el tiempo de acceso ponderado por
distrito para los tres departamentos, en paneles independientes dado que
no son territorios contiguos (un único mapa nacional dejaría la mayor
parte del área en blanco). La escala de color es compartida entre los tres
paneles para permitir la comparación directa: se aprecia visualmente cómo
la severidad de la brecha en Loreto (paneles más oscuros, especialmente en
el extremo este) supera sistemáticamente a la de Ica y Cusco.

\begin{figure}[H]
  \centering
  \includegraphics[width=\textwidth]{figures/gap_map_district_access_time.pdf}
  \caption{Tiempo de acceso ponderado por población, por distrito, para
    Ica, Cusco y Loreto (escala de color compartida). Los distritos en gris
    no tuvieron una coincidencia de nombre exitosa entre la fuente censal
    (INEI, 1999/2002) y el shapefile distrital vigente (IGN).}
  \label{fig:gap-map}
\end{figure}

% =============================================================================
\section{Limitaciones Metodológicas}
\label{sec:limitaciones}
% =============================================================================

Se documentan explícitamente las limitaciones más relevantes para la
interpretación correcta de los resultados:

\begin{enumerate}
  \item \textbf{Resolución temporal censal.} La fuente de centros poblados
    finalmente utilizada (INEI, vía INGEMMET/GEOCATMIN) corresponde al
    censo de \textbf{1999/2002} -- más de dos décadas de antigüedad al
    momento de este estudio. La distribución poblacional real puede haber
    cambiado sustancialmente, particularmente en zonas de crecimiento
    urbano acelerado o de migración. Los pesos poblacionales usados para
    ponderar todos los indicadores heredan esta limitación.
  \item \textbf{100\,\% de los tiempos de esta ejecución son estimaciones
    Haversine, no rutas de red reales.} Como se explicó en la
    Sección~\ref{sec:limitaciones}, sin un contenedor OSRM activo, y dado
    que el muestreo estratificado dispersa la demanda por áreas
    demasiado grandes para la API pública de Overpass, el 100\,\% de los
    pares origen-destino de esta ejecución recurrieron al
    \textit{fallback} de distancia en línea recta $\times$ 1.4. Esto es
    metodológicamente razonable como cota aproximada, pero
    \textbf{subestima sistemáticamente el tiempo real} en zonas de
    orografía compleja (Cusco) o sin red vial continua (Loreto), donde el
    factor de desvío real respecto a la línea recta es probablemente mayor
    a 1.4. Los resultados deben interpretarse como una \emph{cota inferior
    conservadora} de la severidad real de la brecha, no como tiempos de
    viaje reales verificados.
  \item \textbf{Transporte fluvial no modelado.} En Loreto, buena parte de
    la movilidad real ocurre por vía fluvial (ríos navegables), un modo de
    transporte completamente ausente tanto de OpenStreetMap como de
    cualquier extracto OSRM estándar. El modelo asume implícitamente
    transporte terrestre, lo que probablemente \textbf{sobrestima} el
    tiempo de acceso en distritos ribereños con conexión fluvial regular,
    y subestima en aquellos genuinamente aislados por tierra y por río.
  \item \textbf{Muestreo de demanda.} Por razones de viabilidad
    computacional, los indicadores se calculan sobre una muestra
    estratificada del 57.9\,\% del universo de centros poblados (5{,}000
    de 8{,}632), no sobre el universo completo. Aunque el muestreo es
    ponderado por población para minimizar el sesgo, los resultados a
    nivel de distritos con pocos centros poblados (p.\,ej. rankings de
    Tabla~\ref{tab:critical-districts} con $n < 10$ puntos) tienen mayor
    varianza muestral.
  \item \textbf{Truncamiento de datos en La Convención (Cusco).} El
    servicio ArcGIS REST de origen limita cada consulta a 1{,}000
    registros sin soporte de paginación por \textit{offset}; la provincia
    de La Convención, la más extensa y poblada en centros poblados
    pequeños, excedió ese límite incluso paginando por provincia,
    quedando con una cobertura parcial documentada en los logs de
    adquisición.
  \item \textbf{Coincidencia de nombres entre fuentes de distinto vintage.}
    La validación espacial de centros poblados contra su distrito
    declarado (y el mapa de la Figura~\ref{fig:gap-map}) depende de
    coincidencias de nombre entre la nomenclatura distrital del censo
    1999/2002 y el shapefile distrital vigente del IGN. Un número no
    trivial de distritos (373 de 2{,}271 registros de Loreto) no encontró
    coincidencia exacta, probablemente por creaciones, fusiones o cambios
    de nombre de distritos en las últimas dos décadas; también se observó
    al menos un artefacto de codificación de caracteres en un nombre de
    distrito de la fuente censal (``Kos\dh ipata'' en lugar de
    ``Kosñipata''), sin impacto en los cálculos numéricos.
  \item \textbf{Proxy de ruralidad como segundo indicador regional.} El
    índice de ruralidad usado en el cruce analítico se deriva de la propia
    fuente de centros poblados (\texttt{CLASIF02}), por lo que no es
    plenamente independiente del proceso de muestreo de demanda ni de
    variables potencialmente correlacionadas con la calidad de los datos
    de coordenadas.
\end{enumerate}

% =============================================================================
\section{Recomendaciones de Política Pública}
% =============================================================================

A partir de la evidencia presentada, se plantean las siguientes
recomendaciones, priorizadas según el ranking de brechas críticas
(Tabla~\ref{tab:critical-districts}):

\begin{enumerate}
  \item \textbf{Priorizar la frontera amazónica de Loreto para
    intervenciones no convencionales.} Los distritos más críticos del
    ranking (Yavarí, Ramón Castilla, Morona; todos en la provincia
    Mariscal Ramón Castilla o Alto Amazonas) superan los 500 minutos de
    acceso ponderado. Dada la limitación de transporte fluvial no
    modelado, se recomienda complementar este análisis con datos de rutas
    fluviales antes de dimensionar cualquier intervención, y evaluar
    soluciones no viales: telemedicina de emergencia, postas con mayor
    capacidad resolutiva \textit{in situ}, o transporte aéreo/fluvial de
    emergencia subsidiado, antes que infraestructura vial (probablemente
    inviable en el corto plazo en esta geografía).
  \item \textbf{Cusco: cerrar la brecha de La Convención y el corredor sur
    andino.} Los distritos críticos de Cusco (Camanti, Marcapata,
    Kosñipata, Condoroma, Mollepata) están dispersos en varias provincias
    de orografía compleja, no concentrados en una sola. Se recomienda una
    intervención de mejora de vías de penetración (afirmado, puentes) en
    los corredores que conectan estos distritos con sus capitales
    provinciales, donde ya existen establecimientos resolutivos --el
    Gini relativamente bajo de Cusco (0.510) sugiere que el problema es
    más de \emph{cobertura vial dispersa} que de polarización extrema.
  \item \textbf{Ica: intervención focalizada, no estructural.} El Gini más
    alto de los tres departamentos (0.768) combinado con el menor tiempo
    promedio indica que la solución no es una reforma estructural de la
    red, sino intervenciones puntuales en los distritos identificados
    (Marcona, San Pedro de Huacarpana, Huancano) -- probablemente
    ampliaciones de la cartera resolutiva de un establecimiento cercano ya
    existente, dado que la red vial de la región ya es comparativamente
    buena.
  \item \textbf{Focalizar recursos por población, no por número de
    distritos.} La diferencia sistemática entre la media simple y la
    ponderada (Tabla~\ref{tab:access-department}) implica que priorizar
    distritos por conteo simple de unidades geográficas mal servidas
    sobre-representaría a distritos de baja densidad poblacional. Se
    recomienda que cualquier fórmula de asignación presupuestal para
    infraestructura de salud incorpore explícitamente la ponderación
    poblacional aplicada en este estudio.
  \item \textbf{Cerrar la brecha urbano-rural con estrategias
    diferenciadas.} Dada la magnitud de la diferencia (120.6 vs.\ 30.2
    min, $p = 0.0002$), las intervenciones dirigidas a población rural
    dispersa (postas móviles, brigadas periódicas, telemedicina) tienen
    probablemente mayor costo-efectividad marginal que continuar
    invirtiendo en la red urbana ya relativamente bien servida.
  \item \textbf{Actualizar la fuente de centros poblados.} Dada la
    limitación central de vintage censal (Sección~\ref{sec:limitaciones}),
    se recomienda que cualquier decisión de inversión concreta valide los
    hallazgos de este estudio contra el Censo de Población y Vivienda más
    reciente disponible, y contra un extracto OSRM local (en vez del
    \textit{fallback} Haversine) antes de comprometer recursos.
\end{enumerate}

% =============================================================================
\section{Conclusiones}
% =============================================================================

Este estudio construyó un pipeline reproducible de accesibilidad
geoespacial a salud resolutiva para Ica, Cusco y Loreto, documentando de
forma transparente cada decisión técnica --incluyendo las sustituciones de
fuente de datos y las limitaciones del motor de ruteo-- en el propio código
y en este reporte. La evidencia cuantitativa confirma tres hallazgos
robustos a través de múltiples métricas independientes: (i) existe una
brecha de accesibilidad severa y muy desigualmente distribuida entre los
tres departamentos y dentro de cada uno de ellos (Gini nacional 0.729);
(ii) la brecha urbano-rural es estadísticamente significativa y de gran
magnitud (Mann-Whitney $p = 0.0002$); y (iii) la ruralidad distrital está
asociada, aunque no de forma determinante, con un peor acceso
($\rho = 0.266$). Las limitaciones documentadas --particularmente la
dependencia del \textit{fallback} Haversine en ausencia de un motor de
ruteo local y el vintage del censo poblacional-- deben leerse como
invitación a refinar, no a descartar, el diagnóstico: la dirección y el
orden de magnitud de las brechas identificadas son robustos incluso bajo
esas limitaciones, y proveen una base cuantitativa defendible para las
recomendaciones de política planteadas.

\end{document}
